\section{Weighted objectives and attainable regions}
\label{sec:a-weighted}

A linear combination of several potential differences need not inherit the
properties of a single difference. Likewise, a linear combination of arc
flows need not inherit the properties of an individual arc. This section
makes both distinctions precise. Unless stated otherwise, the law is
$g_e(x)=\beta_ex|x|$, all input data are rational, and the connected graph is
simple. Potential objectives are $F=c^\top\pi$ with $\one^\top c=0$;
flow objectives are $d^\top x$. Nomination uncertainty means the balanced
box $\nomset=\{b:l\le b\le u,\ \one^\top b=0\}$ with finite bounds.
Resistance choices are independent positive compact intervals, or explicitly
listed nonempty finite sets where specified. No operating bounds filter
these scenarios. Empty nomination boxes are detected first.

\subsection{Tree hardness and exact rational algorithms}

On a tree, let $w$ be the unique vector satisfying $Aw=c$. Conservation
similarly fixes $x$ from $b$, independently of resistance. Consequently
\begin{equation}
 F=\sum_e w_e\beta_ex_e|x_e|,
 \qquad
 \beta_e^{\rm opt}=\begin{cases}
 U_e,&w_ex_e|x_e|\ge0,\\
 L_e,&w_ex_e|x_e|<0,
 \end{cases}
 \label{eq:a-weight-tree}
\end{equation}
where $L_e,U_e$ are the attained resistance extrema. Both choices belong
to a finite set as well as to its interval hull. Each $w_e$ is an objective
cut sum, and each $x_e$ is a nomination cut sum.

We use the following elementary form of the classical rational quadratic
programming certificate theorem. Its general origin is
\citet{vavasis1990-qp-np}; see also
\citet[Section 2.2, Theorem 3]{delpia2017-miqp-np}.
\begin{lemma}[Rational quadratic maxima]
\label{lem:a-weight-qp}
A rational quadratic polynomial on a nonempty bounded rational polytope
has a rational maximizer of polynomial encoding length. In fixed dimension
$d$, its maximum and a rational maximizer can be found in $N^{O(d)}$ bit
time, where $N$ is the total input length. Singular stationary systems
require no nonsingularity assumption.
\end{lemma}
\begin{proof}
Parameterize the rational affine hull first. Write the quadratic as
$Q(z)=\tfrac12z^\top Hz+a^\top z+a_0$ with $H$ symmetric.
Let a maximum $z^*$ lie in the relative interior of its minimal face,
whose affine hull is $Ez=f$ and whose tangent space has rational basis
columns $T$. Stationarity is the linear system
\[
 Ez=f,\qquad T^\top(Hz+a)=0.
\]
If $z,z'$ satisfy it, then $v=z'-z$ is tangent,
$v^\top(Hz+a)=0$, and subtraction of the stationary equations gives
$v^\top Hv=0$. Thus $Q(z')=Q(z)$. The nonempty intersection of this
affine stationary set with the polytope is bounded; any of its vertices
is rational, has polynomial encoding length by determinant bounds, and
has value $Q(z^*)$. This proves the certificate assertion in any dimension.
For the algorithm, enumerate subsets of at most $d$ independent tight
inequalities, including the empty subset. They include equations for the
affine hull of every face. Solve each stationary system and test its
intersection with the polytope by rational linear programming, retaining
one rational point when feasible. Every retained point is feasible, and
the list includes a global maximum. Compare their rational values.
This handles positive-dimensional stationary sets as well as isolated
stationary points and vertices, using $N^{O(d)}$ operations of polynomial
bit length.
\end{proof}

\begin{theorem}[Arbitrary support on trees]
\label{thm:a-weight-tree-hard}
For quadratic trees with a balanced nomination box, arbitrary rational
zero-sum $c$, and independent finite or interval resistance sets, deciding
whether $\max F\ge T$ is $\NP$-complete. Deciding whether $F\le J$ for every
scenario is $\coNP$-complete. Hardness already holds with fixed resistances,
maximum degree three, and nonzero objective coefficients in $\{-1,1\}$.
Approximating the maximum to absolute error at most $1/8$, or returning
a feasible rational nomination within $1/8$ of it, is $\NP$-hard on
that family. These are absolute-error
statements on an unnormalized class, not strong or relative-error hardness.
\end{theorem}
\begin{proof}
Take positive Subset Sum integers $w_1,\ldots,w_n,K$ with
$0<K<\sum_iw_i$. The omitted trivial targets can be decided in advance
and sent to fixed yes or no instances. Attach a leaf $\ell_i$ to vertex
$v_i$ of the path $v_1,\ldots,v_n$. Give backbone edges resistance one
and edge $v_i\ell_i$ resistance $1/w_i$. Inject $K$ at $v_1$, put zero
nomination at other backbone vertices, and withdraw
$y_i\in[0,w_i]$ at $\ell_i$. Balance is exactly $\sum_i y_i=K$.
All leaf flows are $y_i$, and
\[
 F=\sum_i(\pi_{v_i}-\pi_{\ell_i})
   =\sum_i y_i^2/w_i\le K,
\]
with equality exactly when all $y_i$ are endpoints. This is the classical
convex knapsack endpoint mechanism of
\citet[Proposition 1]{levi2014-convex-knapsack}, here realized with the
stated passive-network restrictions.

A convex maximum over this capped simplex occurs at a vertex. At most
one coordinate is strictly between its bounds; it is an integer because
$K,w_i$ are integers. In a no instance it is some $r\in\{1,\ldots,w_i-1\}$,
and the deficit is
$r-r^2/w_i\ge(w_i-1)/w_i\ge1/2$.
Thus yes maxima equal $K$ and no maxima are at most $K-1/2$.
The robust limit $J=K-1/4$ gives the complementary hardness.
The same gap rules out the asserted additive outputs: a returned rational
nomination can be evaluated exactly. More constructively, with
$m_i=\min\{y_i,w_i-y_i\}$,
$K-F\ge\tfrac12\sum_i m_i$. If $K-F<1/4$, nearest-endpoint rounding
changes the sum by less than $1/2$, so the resulting integer sum equals
$K$. Multiplication of all resistances by $\prod_iw_i$ gives positive
integer resistances of polynomial binary length if desired; the unscaled
gap already proves the approximation assertion.

For membership on the full tree class, eliminate resistance by
\cref{eq:a-weight-tree}. On every closed flow-sign cell of $\nomset$,
the resulting objective is a rational quadratic with polynomial-size
data. A global maximizer lies in one such cell; by
\cref{lem:a-weight-qp} it has a polynomial-bit rational maximizing
nomination. Together with the selected allowed endpoints this is a
certificate. Cut sums and path integration recover rational $x,\pi$ and
verify either $F\ge T$ or the strict violation $F>J$. Existence of this
certificate does not require enumerating the exponentially many cells.
\end{proof}

The comb family also has an exact pseudopolynomial algorithm. For each
possible nonendpoint coordinate $i$, run Subset Sum reachability on the
other capacities up to $K$. For each reachable $s$ with
$0\le K-s\le w_i$, evaluate $s+(K-s)^2/w_i$. A maximum over these
candidates is exact because every capped-simplex vertex has at most one
nonendpoint coordinate; endpoint vertices are included as well.
Backtracking recovers the rational state. This takes $O(n^2K)$ arithmetic
operations on polynomial-bit numbers, not polynomial time in binary input
length.

\begin{theorem}[Objective support and cut directions on trees]
\label{thm:a-weight-tree-algorithm}
On the tree model of \cref{thm:a-weight-tree-hard}, both extrema and a
rational optimizing nomination, resistance, flow, and normalized potential
can be computed in $N^{O(p)}$ bit time when $p=|\supp c|$ is fixed.
A stronger parameter is obtained by contracting all edges with $w_e=0$
and orienting every remaining edge so that $w_e>0$. Let $k$ be the number
of vertices that do not have exactly one incoming and one outgoing edge.
The running time is $N^{O(k)}$, regardless of objective support.
If no edge remains, $F=0$. These bounds are polynomial for each fixed
parameter; no bound of the form $f(p)N^C$ or $f(k)N^C$ is asserted.
\end{theorem}
\begin{proof}
Contracting zero-weight edges is exact: sum the nomination intervals and
objective coefficients in each cluster. Every nomination in a summed
interval can be disaggregated within the original intervals, by filling
from lower bounds greedily. Retained cut flows and all retained terms of
\cref{eq:a-weight-tree} are unchanged; contracted terms are zero.
Rational disaggregation has polynomial bit complexity.

The directed contracted tree has $k$ marked vertices and $k-1$ maximal
paths whose internal vertices each have one incoming and one outgoing
edge. For fixed resistance, smooth the laws to
$\beta_e(x|x|+\rho x)$, $\rho>0$. The nomination gradient $h$ of the
weighted potential objective, defined modulo a constant on the balance
space, satisfies
\[
 h_{\mathrm{tail}(e)}-h_{\mathrm{head}(e)}
   =w_e\beta_e(2|x_e|+\rho)>0.
\]
This follows directly by differentiating the tree cut-sum formula, or
from the adjoint identity in \cref{lem:a-cac-adjoint} with source $c$.
Thus $h$ strictly decreases along each directed path, even if the
magnitudes $w_e$ vary along it. The polyhedral normal cone for a maximum
over the balanced box supplies a common scalar $\lambda$ with
\begin{equation}
 h_v>\lambda\Rightarrow b_v=u_v,\quad
 h_v<\lambda\Rightarrow b_v=l_v,\quad
 l_v<b_v<u_v\Rightarrow h_v=\lambda.
 \label{eq:a-weight-kkt}
\end{equation}
These conditions remain valid on a lower-dimensional nomination domain;
coordinates with equal bounds are already fixed. Along each path the
internal nominations are upper bounds followed by lower bounds, with
at most one free coordinate between them. Enumerate every such pivot,
and every cut without a pivot including both extremes. Leave all marked
nominations free. This gives $N^{O(k)}$ closed faces with at most $2k-1$
free coordinates before balance.

The face family does not depend on $\rho$ or resistance. On the compact
nomination domain, the smoothed objective converges uniformly to $F$:
the added term is $\rho\sum_e w_e\beta_ex_e$, and tree flows are uniformly
bounded. A sequence of smoothed maximizers has a convergent subsequence
on one fixed closed face. Its limit maximizes $F$. For joint optimization,
first fix the resistance of a joint maximizer, then apply this argument.
Thus some joint maximum lies on the same finite family, including for
finite resistance sets.

On a face, parameterize balance and its feasible rational affine hull.
There are at most $2k-1$ variables. Each flow is affine rational in them.
Enumerate the polynomially many cells cut out by the hyperplanes $x_e=0$,
omitting identically zero forms and retaining closed cells; boundaries
cause no problem because a zero-flow term vanishes with either endpoint
choice. The resistance rule \cref{eq:a-weight-tree} leaves a rational
quadratic on each cell. Apply \cref{lem:a-weight-qp}, compare all rational
values, and recover the endpoint resistances and original nominations.
Cut sums and path integration give the whole rational physical state.
Minimization follows by replacing $c$ by $-c$.

For the support bound, every leaf of the contracted tree has nonzero
objective coefficient, so there are at most $p$ leaves and at most $p-2$
branching vertices. A marked degree-two vertex has both currents entering
or both leaving and hence nonzero coefficient. Therefore
$k\le2p-2$. Contraction cannot increase the number of nonzero objective
clusters. This proves the support assertion.
\end{proof}

On a path, $w_e$ is the cumulative sum of $c$ on one side of the cut.
After zero-weight contraction, $t$ changes of sign give $k=t+2$.
In particular, one common sign gives a polynomial algorithm even with
unbounded support. This parameter measures cumulative cut directions,
not signs of the individual coefficients. The pairwise tree algorithms
of \citet[Section 5 of the author manuscript]{labbe2020-bookings-in-the-european-gas} precede this
weighted formulation. Fixed-nomination resistance linearity on trees is
also explicit in \citet[Section 4.2]{amann2018-deciding-robust-feasibility-and-infeasibility};
the additional content here is the weighted nomination-face parameter
and its exact rational output guarantee.

\subsection{A single algebraic core at fixed global cycle rank}

\begin{theorem}[Fixed global rank and supplied signed paths]
\label{thm:a-weight-global}
Fix global cycle rank $r$ and objective support $p$. Under independent
positive resistance intervals, joint optimization of $c^\top\pi$ over
the balanced nomination box admits exact algebraic extrema and an
algebraic optimizing scenario in polynomial bit time for fixed $r,p$.
More generally, after repeatedly removing zero-objective leaves and
aggregating their nomination intervals at their neighbors, suppose a
nonempty marked set $K$ contains every vertex of degree different from
two. Suppose also that on each maximal path with internal vertices
outside $K$, the internal coefficients $c_v$ are all nonnegative or all
nonpositive. Returning paths are allowed. With $k=|K|$ and $r$ fixed,
the same conclusion holds, even when $|\supp c|$ grows. A supplied $K$
and its sign condition can be checked directly. A pure cycle must have
at least one marked vertex. Finite resistance sets are excluded.
\end{theorem}
\begin{proof}
Removing a leaf with $c_v=0$ preserves objective values: its nomination
is aggregated at the neighbor, its resistance can be any allowed value,
and its potential has coefficient zero. Repeat; retain the aggregation
data for recovery. Handle $c=0$ directly. Suppression of the remaining
unmarked degree-two vertices produces a connected multigraph with $k$
vertices and
$P=k-1+r'\le k-1+r$ paths, where $r'$ is its actual rank.
For the support specialization, mark all nonzero coefficients and all
vertices of degree at least three. Every remaining leaf has nonzero
coefficient. The degree identity gives at most $2r'+p-2$ branching
vertices, hence
\[
 k\le2r'+2p-2,\qquad P\le3r'+2p-3.
\]
Parallel and returning paths in the suppressed graph are included.

First hold resistance fixed. Choose $v_0\in K$, give each path a sign
$\sigma\in\{1,-1\}$ agreeing with its internal coefficient signs
(either sign for an all-zero path), and smooth the laws as above.
Perturb the \emph{potential objective} to
\[
 c^\top\pi^\rho+
 \delta\sum_{v\notin K}\sigma_{\mathrm{path}(v)}
                    (\pi_v^\rho-\pi_{v_0}^\rho),\qquad \delta>0.
\]
The adjoint has positive differential resistances
$R_e=\beta_e(2|x_e|+\rho)$ and internal source
$c_v+\delta\sigma$, strictly of sign $\sigma$. On an ordered path,
$j_i=(h_{v_i}-h_{v_{i+1}})/R_i$ satisfies
$j_i-j_{i-1}=c_{v_i}+\delta\sigma$.
Thus the currents strictly increase on positive-source paths and
strictly decrease on negative-source paths. The potential $h$ has a
single peak or valley, respectively; at most one current vanishes,
allowing a plateau of two adjacent vertices at the extremum. Every
horizontal level meets at most two internal vertices, also on a
returning path. By \cref{eq:a-weight-kkt}, positive-source paths have
lower/upper/lower endpoint order and negative-source paths have
upper/lower/upper order, apart from at most two free internal coordinates.
Enumerating these patterns, including missing phases and zero or one
pivot, gives $N^{O(k+P)}$ closed faces with at most
\[
 k+2P=3k+2r'-2
\]
free nominations. Marked coordinates are left free.

This family is independent of both perturbations and resistance.
The uniform limiting argument of \cref{lem:a-blk-faces} applies here
without alteration of its analytic assumptions: total possible
nomination bounds all flows; for $0<\rho\le1$ a path estimate bounds
normalized potentials by sums of $\beta_e(B^2+B)$, where
$B=\sum_v\max\{|l_v|,|u_v|\}$. Compactness and uniqueness of the energy
minimizer give uniform state convergence as $\rho\downarrow0$.
The $\delta$ term tends uniformly to zero. A convergent subsequence of
perturbed maximizers lies on one fixed closed face and yields an original
maximum. Fixing the resistance of a joint maximizer then proves
completeness of the same face family for joint optimization. No numerical
perturbation size is selected by the algorithm.

On one face, eliminate balance and use its rational affine hull with
$d\le k+2P$ variables $z$. A spanning-tree flow gives
$x=x^0(z)+Cq$ with all $r'$ circulation coordinates $q$ retained in the
same core. Chord coordinates equal chord flows and can be bounded by $B$;
free nominations inherit their bounds. Enumerate the polynomially many
flow-sign cells in this fixed-dimensional compact core. On each closed
cell write $x_e|x_e|=p_e(z,q)$, a rational quadratic. The $r'$ cycle
equations and the objective have the form
\[
 \sum_e C_{ej}\beta_ep_e(z,q)=0\quad(1\le j\le r'),
 \qquad F=\sum_e w_e\beta_ep_e(z,q),
\]
where $w$ is any rational solution of $Aw=c$, for example supported on
the spanning tree. This is precisely \cref{lem:a-blk-boxlp}, with
fixed core dimension, fixed aggregate count, degree two, and one scalar
interval leaf $\beta_e$ per edge. It supplies the exact optimum and all
optimizing leaves in a common polynomial-size algebraic representation.
The safe bound $|F|\le B^2\|c\|_1\sum_eU_e$ supplies an explicit objective
bound if needed; $B=0$ is immediate. Compare the polynomially many
algebraic candidates. Disaggregate removed nominations greedily in the
winning field, recover removed flows by conservation, choose any allowed
removed resistance, and integrate drops to recover normalized potentials.
All arithmetic stays in that single field. Replacing $c$ by $-c$ proves
minimization, with path signs reversed.
\end{proof}

\begin{corollary}[Fixed nominations]
\label{cor:a-weight-fixed-nom}
At fixed global rank and fixed rational nominations, independent positive
resistance intervals permit exact algebraic optimization of any rational
zero-sum potential objective or any rational linear flow objective,
without a support bound. With fixed rational resistances the corresponding
value can be compared exactly with a rational threshold in polynomial
time. On trees and with fixed nominations, finite resistance choices
also permit exact rational weighted-potential optimization.
\end{corollary}
\begin{proof}
The preceding algebraic construction has only the $r$ circulation
variables in its core when nominations are fixed. A potential objective
is the same aggregate with arbitrary rational $w$; a flow objective is
affine in the core. Singleton resistance intervals give exact evaluation.
The tree conclusion follows directly from \cref{eq:a-weight-tree}.
\end{proof}

These are global-rank algorithms: they never sum values from independent
algebraic fields. They do not by themselves prove a weighted algorithm
with only block rank bounded. Nor does the theorem give rational optimal
scenarios when cycles remain. The finite-set restriction is substantive,
as the next constructions show.

\subsection{Discrete weighted objectives at ranks one and two}

\begin{lemma}[A three-terminal cycle]
\label{lem:a-weight-triangle}
Orient a triangle $0\to1\to2\to0$, with resistances $(1,1,\theta)$,
$\theta>0$, nominations $(2,-3,1)$, and objective coefficients
$c=(-5,12,-7)$. Its flow is $(q+2,q-1,q)$, where
\[
 6q+3-\theta q^2=0,\qquad
 q=\frac{-6}{6+\sqrt{36+12\theta}}\in(-1/2,0).
\]
Its objective is
\begin{equation}
 F=-5(q+2)^2-7(q-1)^2
   =-\frac{105}{4}-12(q+1/4)^2.
 \label{eq:a-weight-triangle}
\end{equation}
The maximum is $-105/4$ at $\theta=24$, $q=-1/4$.
At $\theta=16/3$ and $144$, respectively, $q=-3/8$ and $-1/8$;
both endpoint values are $-423/16$, below the interior value by $3/16$.
\end{lemma}
\begin{proof}
The displayed flows satisfy conservation. On $(-1/2,0)$ their signs are
$(+,-,-)$, so the cycle equation is
$(q+2)^2-(q-1)^2-\theta q^2=0$.
Its values at $-1/2$ and $0$ have opposite signs. Strict monotonicity
of the sum of signed edge laws gives its unique physical root.
Substitution and completion of the square prove every stated value.
\end{proof}

\begin{theorem}[One-cycle weighted-potential hardness]
\label{thm:a-weight-cycle-hard}
With fixed nominations, independent explicitly listed resistance choices,
and one simple cycle, existential weighted-potential threshold attainment
is $\NP$-complete and robust upper bounds are $\coNP$-complete.
Hardness holds with two-point choices, only three nonzero nominations
$(2,-3,1)$, and only three nonzero objective coefficients $(-5,12,-7)$.
Positive integer resistances suffice. There is a gap of polynomial binary
length, and absolute-error-one value approximation is $\NP$-hard after
polynomially encoded resistance scaling. No strong-hardness or
relative-error claim follows.
\end{theorem}
\begin{proof}
Take positive Subset Sum integers $a_i,K$, put $S=\sum_i a_i$, and
preprocess $K>S$. Replace the triangle's $2\to0$ edge by $n$ series
edges, with zero internal nominations and objective coefficients, and
independent choices
\[
 \beta_i\in\{12/n,\ 12/n+12a_i/K\}.
\]
Their effective resistance is $\theta=12+(12/K)\sum_i a_i\sigma_i$.
Thus \cref{eq:a-weight-triangle} reaches its peak exactly for a target
subset. The graph remains a simple cycle; positive target instances with
one item cause no difficulty. Fixed trivial yes/no outputs use
$\theta=24$ or $12$, with the same peak threshold.

In a no instance $|\theta-24|\ge12/K$. Subtracting the equation at
$q_0=-1/4$ gives
\[
 |q+1/4|=\frac{|\theta-24|}{16(6+\theta(1/4-q))}
 \ge\frac1{20K+12S}.
\]
Here $-1/2<q<0$ and $\theta\le12+12S/K$ bound the denominator by
$240+144S/K$. Hence the no maximum is at most
\[
 -105/4-\Delta,\qquad
 \Delta=\frac{12}{(20K+12S)^2}
       =\frac3{4(5K+3S)^2}.
\]
A strict violation of $J=-105/4-\Delta/2$ exists exactly in the yes
case, and value error below $\Delta/4$ separates the cases. A discrete
scenario within $\Delta/4$ of the maximum also reveals a target subset
in the yes case, checkable with integer arithmetic.

Scaling every resistance by $4nK$ produces first-edge resistances
$4nK$ and path choices $\{48K,48K+48na_i\}$, all positive integers.
Flows are unchanged and the objective and gap scale by $4nK$.
A further factor $(5K+3S)^2$ makes the gap $3nK\ge3$; an absolute-error-one
estimate separates maxima by comparison with the scaled peak minus
$3/2$. Both scalings have polynomial encoding length.

For membership, after guessing one allowed resistance per edge write
$x_e=q+d_e$ in a consistent cycle orientation. Sort the rational
breakpoints $-d_e$ of the strictly increasing cycle equation
$\sum_e\beta_e(q+d_e)|q+d_e|=0$. On each interval it is a rational
quadratic. Rational breakpoint comparisons locate its unique root;
a linear piece or breakpoint root is handled directly. The root has
degree at most two, and all potentials and the weighted objective lie
in that same quadratic field. Exact comparison verifies either a weak
threshold or a strict violation in polynomial bit time. This proves both
membership statements. The interval counterpart is polynomial by
\cref{cor:a-weight-fixed-nom}.
\end{proof}

\begin{lemma}[A rank-two weighted-flow peak]
\label{lem:a-weight-theta}
Use edges $0\to2,2\to1,0\to3,3\to1,2\to3$, with resistances
$(1,1,1,2,\theta)$ and nominations $(4,-4,3,-3)$ at $(0,1,2,3)$.
Writing $a=x_{02}$ and $q=x_{23}$, the unique physical state has
\[
 \begin{split}
 a(q)&=q+9-2\sqrt{2q+18},\\
 x_{21}&=a+3-q,\quad x_{03}=4-a,\quad x_{31}=1-a+q,\\
 \theta q^2&=16-8a(q),\qquad 0<q<3.
 \end{split}
\]
All five flows are positive. With $u=\sqrt{2q+18}$,
\begin{equation}
 F=-9a+5q=-\frac92-2(u-9/2)^2.
 \label{eq:a-weight-theta}
\end{equation}
The peak $-9/2$ occurs at $a=q=9/8$, $\theta_*=448/81$.
At $\theta_L=1792/5329$ the pair $(q,a)=(73/32,57/32)$;
at $\theta_U=12032$ it is $(1/32,17/32)$.
Both endpoint values equal $-37/8$, giving interior advantage $1/8$.
\end{lemma}
\begin{proof}
Balance gives the four outer flows. Equal outer-path drops require
$a^2+(a+3-q)^2=(4-a)^2+2(1-a+q)^2$, and the stated branch solves this.
On $[0,3]$, $a$ increases from $9-6\sqrt2>0$ to
$12-4\sqrt6<4$; also
$a+3-q=12-2\sqrt{2q+18}>0$ and
$1-a+q=2\sqrt{2q+18}-8>0$.
The cross drop is $(4-a)^2-a^2=16-8a$.
The function $\theta q^2-16+8a(q)$ is strictly increasing on $[0,3]$,
negative at zero and positive at three. Its unique root gives positive
flows satisfying both cycle equations, hence the passive state.
Completing the square proves the objective identity; substitution proves
the three rational states.
\end{proof}

\begin{theorem}[Rank-two weighted-flow and total-flow hardness]
\label{thm:a-weight-flow-hard}
At fixed nominations and global rank two, existential linear-flow
threshold attainment with independent finite resistance sets is
$\NP$-complete and robust upper bounds are $\coNP$-complete.
Hardness holds on maximum-degree-three graphs with two-point choices,
the four fixed nominations in \cref{lem:a-weight-theta}, and two fixed
objective coefficients $(-9,5)$, or $(9,5)$ with a shifted threshold.
It also holds for unit-cost total absolute flow on a directed acyclic
graph with those same small nominations. Positive integer resistances
suffice. These reductions have polynomial-bit gaps, obstructing
polynomial dependence on accuracy bits unless $\mathrm P=\NP$.
Absolute-error-one hardness follows separately by scaling nominations;
common resistance scaling alone does not enlarge a flow gap.
\end{theorem}
\begin{proof}
Take positive Subset Sum data $a_i,K$. In the theta of
\cref{lem:a-weight-theta}, replace the cross edge by an $n$-edge path with choices
\[
 \left\{\frac{224}{81n},
 \frac{224}{81n}+\frac{224a_i}{81K}\right\}.
\]
Its total resistance is
$\theta=(224/81)(1+\sum_i a_i\sigma_i/K)$, attaining $448/81$
exactly for a target subset. Use any cross-path edge for the $5q$ term.
The global rank is two, degree at most three, and all subdivision
nominations are zero. Preprocess $K>S=\sum_i a_i$ with fixed yes/no
theta instances $448/81$ and $224/81$.

For the gap put $H=23K+15S$. In a no instance,
$|\theta-\theta_*|\ge224/(81K)>2/K$ and
$\theta\le3(1+S/K)$. Subtract the cross equations at $q$ and $q_*=9/8$:
\[
 (\theta-\theta_*)q_*^2
 =-(q-q_*)\left(\theta(q+q_*)+
       8\frac{a(q)-a(q_*)}{q-q_*}\right).
\]
The quotient is in $(0,1)$ since $0<a'(q)<1$; also $q+q_*<5$ and
$q_*^2>1$. Hence $|q-q_*|>2/H$.
Since $u+9/2<10$, $|u-9/2|>2/(5H)$.
The no maximum is therefore at most $-9/2-\Delta$, where
$\Delta=1/(4H^2)$. Use the strict robust limit $-9/2-\Delta/2$ or
value error below $\Delta/4$. A near-optimal discrete scenario similarly
reveals a target subset in the yes case.

Scaling all resistances by $81nK$ multiplies the four outer coefficients
$(1,1,1,2)$ by $81nK$ and gives cross choices
$\{224K,224K+224na_i\}$. It changes no flows, values, or gaps.
For absolute-error-one hardness, instead multiply all nominations by
$M=16H^2$ after this integer conversion. Homogeneity
(\cref{lem:a-pre-homogeneity}) scales flows and $F$ by $M$, so the gap
is at least four and comparison with the scaled peak minus two suffices.
The nominations in this last assertion are no longer fixed small numbers.

For positive coefficients, conservation at $0$ gives
$9x_{03}+5x_{23}=F+36$, with target $63/2$ and the same gap.
For unit coefficients on \emph{every} edge, subdivide $0\to3$ into
$9n+1$ edges of total resistance one. Keep the other three outer edges.
Use $5n$ cross-path edges: $4n$ fixed edges of resistance $28/(81n)$
and $n$ uncertain edges with choices
\[
 \left\{\frac{112}{81n},
 \frac{112}{81n}+\frac{224a_i}{81K}\right\}.
\]
The effective cross resistance is unchanged. All flows are positive and
the orientation is acyclic (order $0,2,3,1$, inserting path vertices).
Thus total absolute flow equals total oriented flow, namely
\[
 (9n+1)(4-a)+a+(a+3-q)+(1-a+q)+5nq
   =nF+36n+8.
\]
Its peak is $(63n+16)/2$ and its gap at least $n\Delta$.
There are $14n+4$ edges and $14n+3$ vertices, preserving rank two and
maximum degree three. Common resistance scaling by
$81nK(9n+1)$ makes all data integral: the long outer-path edges have
resistance $81nK$, fixed cross edges $28K(9n+1)$, and uncertain options
$112K(9n+1)$ and $112K(9n+1)+224n(9n+1)a_i$.
No constant-gap claim with the small nominations is inferred.

For membership, guess the finite choices and use
\cref{cor:a-weight-fixed-nom} with singleton intervals: its fixed-rank
algebraic evaluation verifies weak attainment or strict violation exactly.
For total absolute flow, flow-sign cells make the objective affine in
the core as well. The constructions therefore prove the stated complete
problems, not just hardness.
\end{proof}

\subsection{Cactus flow boxes and the difference between scenarios and values}

\begin{theorem}[Flow region and exact optimizing resistance scenario]
\label{thm:a-weight-cactus-region}
Fix rational balanced nominations on a cactus. With independent resistance
intervals, its attainable flow region is an affine image of independent
scalar circulation intervals. With finite resistance sets, the convex hull
of the attainable flow region is that same image for the interval hulls.
For any rational linear flow objective, a rational endpoint resistance
scenario attaining the exact optimum can be found in polynomial bit time,
even with arbitrarily many cycles. The value admits an additive estimate
in time polynomial in input size and accuracy bits. Exact comparison of
that scalar value with a rational threshold is nevertheless $\SRS$-hard,
already with fixed positive integer resistances and unit source/sink
nominations. No $\NP$-hardness claim is made for this arithmetic barrier.
\end{theorem}
\begin{proof}
By \cref{lem:a-pre-block,lem:a-pre-cycle}, bridge flows are fixed and a
consistently oriented cycle has flows $q+d_e$, with rational fixed offsets.
Its physical circulation solves
$H_\beta(q)=\sum_e\beta_e(q+d_e)|q+d_e|=0$.
Define
\[
 H_-(q)=\sum_e\min\{L_e(q+d_e)|q+d_e|,U_e(q+d_e)|q+d_e|\},
\]
\[
 H_+(q)=\sum_e\max\{L_e(q+d_e)|q+d_e|,U_e(q+d_e)|q+d_e|\}.
\]
Both are continuous and strictly increasing, with opposite infinite
limits. Let $q^+$ be the root of $H_-$ and $q^-$ the root of $H_+$.
Since $H_-\le H_\beta\le H_+$, all attainable circulations lie in
$[q^-,q^+]$. At either root choose each endpoint attaining its summand;
this gives an actual endpoint scenario realizing that root, including
for finite sets. Sorting the rational breakpoints $-d_e$ makes both
equations piecewise quadratic. Rational sign tests locate the appropriate
piece; its root, including linear or breakpoint cases, has a polynomial-size
algebraic encoding of degree at most two. Comparing it to the breakpoints
recovers the realizing resistance endpoints exactly in polynomial time.
This scalar root principle is consistent with the cycle-interval
linearization in \citet[Lemma 4.10]{amann2018-deciding-robust-feasibility-and-infeasibility}.

For intervals, the continuous image of the connected compact resistance
box is the whole interval $[q^-,q^+]$. Different blocks have fixed
effective nominations and disjoint resistance variables. Thus, for a
rational particular flow $x^0$ and signed cycle matrix $Z$,
\begin{equation}
 \mathcal X=\{x^0+Zq:q_C\in[q_C^-,q_C^+]\text{ for all cycles }C\}.
 \label{eq:a-weight-flow-box}
\end{equation}
Cycle supports are disjoint, so this is a possibly degenerate
parallelotope. For finite sets the attainable circulation set is a
Cartesian product; the convex hull of a finite Cartesian product is the
product of its convex hulls. Each scalar hull is the same interval,
proving the second assertion.

A linear flow objective is $D+\sum_C A_Cq_C$, with rational $D,A_C$.
Choose the endpoint scenario for $q_C^+$ when $A_C>0$, for $q_C^-$ when
$A_C<0$, and any allowed setting when $A_C=0$. This requires no comparison
of sums of independent radicals. Approximate each selected root within
$\eps/[2(1+\sum_C|A_C|)]$ and sum to obtain value error below $\eps$.
Quadratic root enclosure takes polynomial time in the input and accuracy
bits. Separate quadratic encodings can also represent the physical flow;
a common field containing all cycle roots is not required.

For the arithmetic obstruction take a $\SRS$ instance
$\sum_i\sqrt{a_i}\le K$ with positive integer data. Remove $a_i=1$ terms and subtract their count
from $K$. If no terms remain or the adjusted $K\le0$, decide directly
and use fixed outputs below. Otherwise build a chain of separate triangles
joined by resistance-one bridges with unit injection at the first and
unit withdrawal at the last terminal. In triangle $i$, a direct edge of
resistance $a_i>1$ is parallel to a two-edge path of resistances $1/2,1/2$.
The direct flow is $x_i=1/(1+\sqrt{a_i})$. Give it objective coefficient
$a_i-1$ and give all other edges coefficient zero. With $m$ triangles,
\[
 F=\sum_i(a_i-1)x_i=\sum_i\sqrt{a_i}-m.
\]
Thus $F\le K-m$ is exactly the source comparison, preserving equality.
Doubling every resistance makes all data positive integers and leaves
flows unchanged. Distinct triangles joined by bridges give a simple
maximum-degree-three cactus. Fixed yes/no outputs are the triangle with
direct resistance $8$, alternate resistances $1,1$, and direct-edge
objective coefficient $3$: its value is $1$, so thresholds $1$ and $0$
serve. The $\SRS$ convention and its unresolved complexity are as in
\cref{subsec:a-pre-computation}. Here the scenario is unique, so scenario
selection is trivial while exact scalar comparison still contains $\SRS$.
\end{proof}

For rank at most one, the objective in the proof has at most one quadratic
root, so even exact scalar comparison is polynomial. Thus rank two is
the first possible discrete weighted-flow hardness rank in
\cref{thm:a-weight-flow-hard}, while arbitrary cactus rank creates a
separate arithmetic issue. Rational membership in \cref{eq:a-weight-flow-box}
is decided by conservation and separate quadratic interval comparisons.
For finite sets this tests convex-hull membership, not finite attainability,
which is already hard on one cycle by \cref{thm:a-bound-realization}.
Every continuous convex function on the flow box has a maximizing vertex,
since its value at a convex combination is bounded by the largest vertex
value. Each vertex has an endpoint resistance realization. This structural
fact is not an efficient algorithm for arbitrary convex maximization.

\subsection{The objective determines the universal resistance-hull class}

Call a graph a \emph{universal hull graph} for an objective class if, for
every fixed balanced nomination and every objective in that class,
independent nonempty compact positive resistance sets have the same
maximum and minimum as their interval hulls. The sets have attained
endpoints by compactness; no representation of their interiors is needed
for this definition. Separate monotonicity means monotonicity of every
one-coordinate section after all other coordinates are fixed. The
direction may depend on those fixed coordinates and on the nomination.

\begin{theorem}[Universal hull hierarchy]
\label{thm:a-weight-hierarchy}
Under quadratic laws, finite connected simple graphs have the following
universal hull classes:
\[
\begin{array}{c|c}
\text{objective class}&\text{graph class}\\ \hline
\text{all zero-sum linear potential objectives}&\text{trees}\\
\text{all individual potential differences}&\text{cacti}\\
\text{all linear flow objectives}&\text{cacti}\\
\text{all individual arc flows}&\text{series--parallel graphs}.
\end{array}
\]
In each row the same graph class is characterized by universal separate
monotonicity over positive resistance boxes. Each negative direction has
a witness with one varying resistance, rational nominations and objective
coefficients, and strictly positive rational resistances, all of polynomial
encoding length, given the relevant cycle or theta subdivision
(or the $K_4$ subdivision in the last row).
\end{theorem}
\begin{proof}
The final row is \cref{thm:a-bound-characterization}. We prove the other
three, including the restoration bounds needed on graphs larger than
the basic witnesses.

\emph{Positive directions.}
On trees, fixed nominations fix every flow, and
\cref{eq:a-weight-tree} is affine in resistance. For pairwise potentials
on cacti, the cactus part of \cref{thm:a-bound-envelope-structural}
gives the required one-coordinate monotonicity: cactus terminal adjoint currents have fixed
signs along their unique block route. To spell out why the sign of the
law change stays fixed, for an affine one-edge parameter
$g_e(x,\alpha)=g_e^0(x)+\alpha f_e(x)$, a zero of
$f_e(x_e(\alpha))$ makes the entire state at that parameter satisfy the
laws at every parameter value. Uniqueness makes the state constant.
Otherwise continuity fixes the sign of $f_e(x_e)$, and the finite-secant
identity fixes the direction of the terminal objective. This also proves
the positive direction for the admissible continuous affine edge-law
families of \cref{thm:a-bound-envelope-structural}, including several
independent coordinates on one edge.

For linear flows on a cactus, block independence makes the objective
$D+\sum_C A_Cq_C$. Vary only one resistance $\beta_e$ in one cycle.
At $q=-d_e$, the strictly increasing cycle function $H_\beta(q)$ is
independent of $\beta_e$, so the physical sign of $q+d_e$ is fixed as
$\beta_e$ varies. If it is zero once, the root is constant. Otherwise
increasing $\beta_e$ changes the old root's equation with fixed nonzero
sign, so root comparison makes $q$ monotone. Its affine objective
contribution is monotone as well.

In every positive direction, start at an optimizer over the compact
interval box and move each coordinate to a maximizing endpoint without
reducing the objective. Its value remains the global optimum after every
move, even if later monotonicity directions change. The final corner
belongs to the original product of compact sets. Apply the same argument
to the negative objective for minima. Thus separate monotonicity implies
the stated hull equality.

\emph{Weighted-potential converse.}
On any non-tree choose a simple cycle and three distinct vertices on it.
Represent its three connecting paths as subdivisions of
\cref{lem:a-weight-triangle}. Split each resistance-one side into equal
positive shares. On the uncertain side, reserve fixed total $8/3$ on
all but one edge if subdivided, and put resistance $\theta-8/3$ on the
remaining edge; on a single edge use $\theta$. The three values
$\theta\in\{16/3,24,144\}$ remain positive and reproduce the triangle
because internal nominations vanish. Put the three stated nominations
and objective coefficients on its branch vertices and zero elsewhere.

Let $m=|E|$ and restore every other edge with resistance
\[
 R=12(10000m^2)^3.
\]
The conserved comparison flow with $q=-1/4$ on the selected cycle and
zero elsewhere has energy
$[(7/4)^3+(5/4)^3+\theta(1/4)^3]/3=(468+\theta)/192<4$
at all three settings. The full physical minimizer therefore has every
extra-edge flow at most $u=(12/R)^{1/3}$ in magnitude. Total positive
nomination is three, so every full flow has magnitude at most three.
Restrict the full state to the cycle; its induced nomination $b'$ is
balanced, has $|b'_v|\le6$, and satisfies
$\|b'-b\|_1\le2mu$. The original and induced nominations lie in a common
balanced box with total absolute bound $B\le6m$.
Apply \cref{lem:a-cac-lipschitz} to the two paths of total resistance one
in $F=-5(\pi_0-\pi_1)+7(\pi_1-\pi_2)$. The objective error at each
setting is at most
\[
 (5+7)\,2(6m)\,(2mu)=288m^2u=288/10000<3/64.
\]
The restored interior value exceeds both endpoints by more than
$3/16-2(3/64)=3/32$. This violates both properties.

\emph{Theta subdivisions.}
Every noncactus simple connected graph contains three internally disjoint
paths between two vertices. Indeed, take a block other than an edge or
cycle and a cycle in it. A chord gives such a theta. Otherwise an
outside component has at least two distinct neighbors on the cycle,
since one neighbor would be an articulation; an internal connecting
path and the two cycle arcs give the theta. At most one theta path is
a single edge by simplicity. Designate branch vertices $2,3$, choose
internal vertices $0,1$ on the other two paths, and split the four
outer segments at them. This realizes either four-vertex theta gadget
by positive resistance subdivisions with zero internal nominations.
Keep just one edge on the cross path variable, reserving a fixed total
less than its lower resistance endpoint on all other cross edges.

\emph{Pairwise-potential converse.}
Use the pressure theta already calculated in
\cref{eq:a-bound-theta}, with outer resistance totals
$(1/2,1/6,1/6,1/2)$ and nominations $(4,-4,3,-3)$.
No leaf from that hardness reduction is needed. For cross total
$\tau\in\{23/108,1,71/12\}$ the cross flows are respectively
$3/2,1,1/2$, and
\[
 \pi_1-\pi_0=-2-(q-1)^2/6.
\]
The interior exceeds both endpoints by $1/24$.
For a subdivided cross path reserve fixed total $23/216$ and vary the
remaining edge, preserving positivity. Restore all extra edges with
\[
 R=18(10000m^2)^3.
\]
The comparison flow at $q=1$ has outer flows $(1,3,3,1)$ and energy
$(10+\tau)/3<6$. Extra flows are at most $u=(18/R)^{1/3}$.
The full-flow bound $14$ gives
$|b'_v|\le14\deg_{\rm theta}(v)$ and
$B\le28m_{\rm theta}$ on the selected subgraph, with
$\|b'-b\|_1\le2m_{\rm extra}u$. The terminal path through $2$ has
resistance $2/3$, so \cref{lem:a-cac-lipschitz} bounds the pressure error by
\[
 \frac{224}{3}m_{\rm theta}m_{\rm extra}u
 \le75m^2u=75/10000<1/96.
\]
Thus the strict interior advantage survives.

\emph{Linear-flow converse.}
Use \cref{lem:a-weight-theta} with its three cross resistances
$1792/5329,448/81,12032$. Split outer totals $(1,1,1,2)$ on the theta
segments and, if necessary, reserve fixed cross total $896/5329$.
Choose one edge of segment $0\to2$ for coefficient $-9$ and one
cross-path edge for coefficient $5$. On the isolated theta their flows
are $a,q$. Restore all other edges with
\[
 R=18000(10^9m^2)^3.
\]
The conserved comparison flow $a=q=9/8$ has other path flows
$(3,23/8,1)$ and energy
\[
 \bigl[(9/8)^3+3^3+(23/8)^3+2+\theta(9/8)^3\bigr]/3<6000
\]
at all three settings (at the upper setting it is $91657/16$).
Hence extra flows are at most $u=(18000/R)^{1/3}$.
All full flows are bounded by the total positive nomination seven.
The induced theta nominations obey $|b'_v|\le21$,
$B\le21m$, and $\|b'-b\|_1\le2mu$.
For either selected edge, scalar quadratic monotonicity gives
\[
 \begin{split}
 |x_e(b')-x_e(b)|^2
 &\le\frac2{\beta_e}
   |(\pi_{\rm tail}-\pi_{\rm head})(b')
        -(\pi_{\rm tail}-\pi_{\rm head})(b)|\\
 &\le4B\|b'-b\|_1\le168m^2u,
 \end{split}
\]
using \cref{lem:a-cac-lipschitz} on the single-edge path. This follows
also directly from
$|s|s-|t|t$ having magnitude at least $|s-t|^2/2$.
The resistance cancels even on a subdivided uncertain edge.
The objective error is at most $14\sqrt{168m^2u}<1/32$, since its
square is at most $196\cdot168/10^9<1/1024$.
The restored interior advantage exceeds $1/8-2(1/32)=1/16$.
Alternatively select the $0\to3$ segment with coefficient $9$ and the
cross edge with coefficient $5$. On the isolated theta the objective
is $F+36$; after restoration the same error estimate applies directly
to those selected edges. Thus positive objective coefficients suffice.

All induced nominations above are balanced because the added vertices
have zero original nomination and aggregate extra-edge exchanges cancel.
The restricted physical state satisfies the selected subgraph's laws,
so uniqueness justifies each comparison with its perturbed-nomination
state. All resistance splits and displayed restoring integers have
polynomial encoding length. Each witness has a strict interior value
above both endpoints, and its endpoint-only set has a smaller maximum
than its interval hull. This completes all converses.
\end{proof}

The positive hull statements also hold jointly over any compact
nomination set for which the state is continuous: fix the nomination
of a joint optimizer, then move resistance coordinates to endpoints.
This does not assert separate monotonicity of the nomination-optimized
value. For pairwise potentials, the continuous affine one-edge parameter
version proved above extends beyond quadratic resistances; a parameter
shared by several edges is not covered.

There is a constructive additive endpoint recovery when the scalar sets
have rational attained endpoints. Suppose the applicable interval-hull
algorithm returns a rational feasible nomination and parameter scenario
within $\eps/2$ of its maximum. Keep the nomination fixed and, for each
of $M>0$ parameters, evaluate the two endpoint settings within
$\delta=\eps/(4M)$ and retain the larger estimate. Separate monotonicity
ensures that the better endpoint is no worse than the current value;
selecting by these estimates loses at most $2\delta$. After $2M$
fixed-scenario evaluations the total loss is at most $\eps/2$, and every
parameter belongs to its original compact set. Minima use the smaller
estimate. The polynomial and fixed fractional law algorithms of
\cref{sec:a-laws} supply the needed accuracy-bit evaluations in their
respective scope; this argument presupposes such an evaluation algorithm
and does not assign one to an arbitrary continuous law. For $M=0$ no
rounding is needed. In particular, no exact comparison of a sum of
independent pressure radicals is hidden in endpoint recovery.

\subsection{Convexity of entire attainable state regions}

Normalize one potential to zero throughout this subsection. An attainable
region varies resistance in a positive interval box while fixing balanced
nominations; the output may be the whole flow vector, the whole normalized
potential vector, or their joint vector.

\begin{lemma}[One varying cyclic resistance]
\label{lem:a-weight-deletion}
For any strictly increasing continuous passive edge laws of the form
$\beta_ef(x_e)$ with $f(0)=0$, vary just one resistance on a cyclic edge
$e=(u,v)$. Either the entire state is constant, or its own-edge flow $q$
is a continuous strictly monotone parameter for the attainable state.
In the latter case the terminal drop $\pi_u-\pi_v$ is also an injective
linear coordinate on the potential region. If either the flow region
or the potential region is convex, that region equals the segment between
its endpoint states, and every linear objective has endpoint extrema.
\end{lemma}
\begin{proof}
Delete $e$, leaving a connected graph. Prescribing its flow $q$ gives
remaining nominations $b-q(e_u-e_v)$. Let $P(q)$ be the remaining
network's terminal drop. For $q_2>q_1$, its two conserved flows differ,
and strict monotonicity implies
\[
 0<\sum_{f\ne e}[g_f(x_f(q_2))-g_f(x_f(q_1))]
                      [x_f(q_2)-x_f(q_1)]
   =-(q_2-q_1)[P(q_2)-P(q_1)].
\]
Thus $P$ is continuous strictly decreasing. The full state solves
$P(q)=\beta_ef(q)$, an equation with unique root. Its sign is determined
by $P(0)$ independently of resistance. If $P(0)=0$, then $q=0$ for all
resistances and the remaining state is fixed. Otherwise changing
$\beta_e$ changes the old root's equation by a nonzero term of fixed
sign, so root comparison makes $q(\beta_e)$ strictly monotone.
State continuity gives a continuous bijection onto its interval of flows.
The remaining state is uniquely determined by $q$, and $P(q)$ is strictly
monotone. Therefore each flow vector is uniquely determined by its
coordinate $q$, and each potential vector by its linear coordinate
$\pi_u-\pi_v=P(q)$. If the region is convex, its endpoint chord belongs
to it and already contains every intermediate coordinate value. Uniqueness
above each coordinate forces the entire region to be that chord.
\end{proof}

\begin{theorem}[Universal convex-region classes]
\label{thm:a-weight-convex-regions}
For the quadratic law on a finite connected simple graph, the attainable
flow region is convex for every fixed balanced nomination and every
positive resistance box if and only if the graph is a cactus.
The analogous class for normalized potential regions, and for joint
flow--potential regions, is exactly the trees. Negative witnesses can
have one varying resistance, rational nominations, and strictly positive
rational resistances, all of polynomial encoding length.
The same universal classes hold if every interval is required to
have strictly positive width, but the positive-width extension below is
existential and supplies no polynomial encoding bound on the widths.
\end{theorem}
\begin{proof}
The cactus flow region is the affine box in
\cref{thm:a-weight-cactus-region}. On a tree, flows are fixed and normalized
potentials are linear in resistance; both potential and joint regions
are affine images of the box. On every noncactus, the one-coordinate
linear-flow witness of \cref{thm:a-weight-hierarchy} contradicts the
endpoint conclusion of \cref{lem:a-weight-deletion}, so its flow region
is nonconvex. On every non-tree, the weighted-potential witness contradicts
the corresponding potential conclusion. Its joint region cannot be
convex because its linear projection onto potentials is nonconvex.
In both applications the strict interior advantage excludes the
constant-state exception in the deletion lemma.

For strictly positive widths, take one of these compact nonconvex
regions $S$ and a point $z$ on a chord between two of its states with
$z\notin S$. Compactness gives $\operatorname{dist}(z,S)>0$.
Thicken all fixed resistance coordinates to small positive-width rational
intervals containing their original positive values. The two original
states remain attainable. Uniform continuity on a compact positive
resistance neighborhood implies that all states of a sufficiently small
thickening remain within less than half that distance of $S$.
They still exclude $z$, so the new region remains nonconvex. This applies
to flows, normalized potentials, and joint states. Continuity supplies
existence of rational widths, not their bit complexity.
\end{proof}

\begin{theorem}[Qualitative extension to every nonlinear power]
\label{thm:a-weight-power-regions}
For each fixed real exponent $\gamma>0$, $\gamma\ne1$, the classifications
in \cref{thm:a-weight-convex-regions} remain valid for the common law
$g_e(x)=\beta_e\sgn(x)|x|^\gamma$. Witnesses with rational nominations
and strictly positive rational resistances exist. This is a geometric statement; no uniform
encoding bound, exact algebraic algorithm, or accuracy-bit algorithm is
asserted for this extension. The linear exponent $\gamma=1$ is excluded.
\end{theorem}
\begin{proof}
Strict convexity and coercivity of
$\sum_e\beta_e|x_e|^{\gamma+1}/(\gamma+1)$ give a unique continuous
state. On cacti, the image of each independent compact connected
resistance box in its scalar circulation coordinate is an interval,
proving the flow positive direction. On trees, the fixed flow makes
potentials linear in resistance, proving the other two.

For a flow obstruction use the theta edges of
\cref{lem:a-weight-theta}, now with outer resistances $(1,3,2,2)$ and
nominations $(2,-2,0,0)$. Conservation gives outer flows
$(a,a-q,2-a,2-a+q)$. Near $(a,q)=(1,0)$ all are positive, and their
outer-cycle equation is
\[
 \Phi(a,q)=a^\gamma+3(a-q)^\gamma
               -2(2-a)^\gamma-2(2-a+q)^\gamma=0.
\]
At that point $\Phi_a=8\gamma>0$. Implicit differentiation gives
$a'(0)=5/8$. Moreover the second-derivative numerator divided by
$\gamma(\gamma-1)$ is
\[
 (1-2)(5/8)^2+(3-2)(5/8-1)^2=-1/4,
\]
so $a''(0)=(\gamma-1)/32\ne0$.
The cross drop $P(q)=2(2-a(q))^\gamma-a(q)^\gamma$ equals one at zero
and stays positive for small positive $q$. The resistance
$\beta_{23}(q)=P(q)/q^\gamma$ therefore realizes these states.
The deletion lemma makes this continuous parametrization injective;
choose a compact nontrivial positive $q$ interval on which $a''$ has
strict sign. Its graph is not a segment. If $L(q)$ is the endpoint
secant of $a$, then either $a-L(q)$ or its negative is a linear flow
objective plus a constant with strictly positive interior value and zero
endpoint values. The constant is immaterial.

For potentials take the triangle consisting of path $0\to1\to2$ with
both resistances one and a variable direct edge $0\to2$, with
nominations $(2,1,-3)$. The path flows are $a,a+1$, and the direct flow
is $2-a$. For $0<a<2$ choose
\[
 \beta_{02}(a)=\frac{a^\gamma+(a+1)^\gamma}{(2-a)^\gamma}.
\]
It is strictly positive and strictly increasing. Normalizing $\pi_2=0$
gives $(\pi_0,\pi_1)=(a^\gamma+(a+1)^\gamma,(a+1)^\gamma)$ and
\[
 \frac{d\pi_0}{d\pi_1}
   =1+\left(\frac a{a+1}\right)^{\gamma-1}.
\]
This has nonzero derivative on an interior interval when $\gamma\ne1$.
The potential graph has strict curvature, hence a linear potential
objective (made zero-sum by its coefficient at $2$) with strict interior
advantage over its endpoints. Its potential region is nonconvex by
\cref{lem:a-weight-deletion}; so is its joint region.

Subdivide these witnesses on a theta in any noncactus, or a cycle in any
non-tree, exactly as in \cref{thm:a-weight-hierarchy}. Series resistance
coefficients add for every common exponent; all new internal nominations
are zero. Reserve a small fixed positive cross-path total to keep only
one resistance variable. Restore all other edges with a common resistance
$R$. A conserved comparison flow supported on the selected subgraph has
uniformly bounded energy for the compact interval of varying resistance.
Energy minimality makes every extra flow $O(R^{-1/(\gamma+1)})$ uniformly.
The full physical flow bound is independent of $R$. Restricting to the
selected subgraph therefore induces nomination perturbations tending
uniformly to zero. Uniqueness and continuous dependence of that subgraph's
state show uniform convergence of its selected flows and normalized
potentials. The strict interior objective advantage survives for some
finite $R$. The variable edge is cyclic in the full graph, so the
deletion lemma again turns that advantage into nonconvexity of the
corresponding entire region.

Finally, strict inequalities and continuity allow rational approximations
of the two resistance endpoints and the chosen interior setting, and a
sufficiently large rational $R$. Resistance subdivisions can use positive
rational shares; the displayed nominations are already integers. If the
separating objective initially has real coefficients, approximate these
as well while preserving its strict gap and zero sum when required.
This proves existence with rational data without bounding their sizes.
The positive-width argument of \cref{thm:a-weight-convex-regions} applies
unchanged. At $\gamma=1$ both curvature computations vanish, so neither
converse is asserted by this proof.
\end{proof}

These classifications concern convex sets of full vectors under
\emph{resistance variation}. They differ from convex scalar transfer
characteristics under source variation. In particular,
\citet[Theorem 4, p.~13, and conclusion, p.~21]{wanghasler1997-convexity}
give topological conditions for the latter and propose extension to
other circuit parameters as further work. The separate
1993 nonlinear-tolerance paper
\citep{hasler1993-parameter-tolerances} has not been read in full, so we do
not make an exhaustive priority claim. The curvature, energy, and
convex-graph arguments used here are elementary mechanisms; the
specific universal output classes, restricted converse witnesses, and
quantified quadratic restoration are the statements established here.

A related cactus characterization appears in
\citet[Definitions 1--2 and Theorem 18 of the preprint]{brandenberg2025-differential}:
each linear differential-flow polytope fixes its positive elasticity
vector, while their universal nondegeneracy characterization quantifies
over all positive elasticity vectors and all nomination/capacity bounds.
Those quantifiers and that property differ from resistance-attainable
nonlinear state regions. For an additional contrast, with linear Ohmic
laws \emph{every} graph has the potential-objective hull property.
Indeed, ground one vertex, vary one conductance from $t_0$ to $t$, and
write $L(t)=L_0+(t-t_0)aa^\top$. For any grounded objective vector $c$,
put $U=c^\top L_0^{-1}a$, $V=a^\top L_0^{-1}b$, and
$r=a^\top L_0^{-1}a$. The rank-one inverse identity gives
\[
 c^\top\pi(t)=c^\top\pi(t_0)
        -\frac{(t-t_0)UV}{1+(t-t_0)r}.
\]
Positive definiteness makes the denominator positive for every positive
conductance in the interval, so the derivative has fixed sign $-UV$.
Reciprocal resistance reverses the parameter direction and preserves
monotonicity. Coordinatewise endpoint selection then proves the linear
hull assertion on all connected graphs. This familiar sensitivity
calculation does not assert convexity of their entire state regions.
